We evaluate AVControl on standard benchmarks, demonstrate its extendability across diverse modalities, and analyze training efficiency.

\subsection{Experimental Setup}
\label{sec:setup}
We train all LoRAs on top of LTX-2~\cite{hacohen2026ltx2}, a frozen joint audio-visual foundation model. Each per-modality LoRA is trained independently on a single H100 GPU; full training details are provided in the supplementary (Table~\ref{tab:training_details}).

We use fixed generation parameters across all evaluations: constant seed~42, CFG~1.0, and LoRA strength~1.0. The low guidance scale reflects LTX-2's distilled inference mode; the control signal provided via the parallel canvas supplies sufficient structural context, making high guidance scales unnecessary.

For quantitative evaluation, we adopt the VACE Benchmark~\cite{jiang2025vace} (20~samples each for depth, pose, inpainting, and outpainting). We use the exact same published input videos and control signals as all baselines, ensuring a fair comparison against reported numbers. We report six VBench~\cite{huang2024vbench} metrics: Aesthetic Quality (AQ), Background Consistency (BC), Dynamic Degree (DD), Imaging Quality (IQ), Motion Smoothness (MS), and Subject Consistency (SC).

\subsection{Quantitative Evaluation}
\label{sec:quant}

\begin{table*}[t]
\centering
\caption{VBench evaluation on the VACE Benchmark. All methods use the same inputs. Baseline numbers from~\protect\cite{jiang2025vace}. \textbf{Bold}: best per metric.}
\label{tab:vbench_comparison}
\setlength{\tabcolsep}{4pt}
\begin{tabular}{llccccccc}
\toprule
Task & Method & AQ & BC & DD & IQ & MS & SC & Avg. \\
\midrule
\multirow{5}{*}{Depth} 
& Control-A-Video~\cite{chen2023controlavideo} & 50.6 & 91.7 & \textbf{70.0} & 67.8 & 97.6 & 88.1 & 77.6 \\
& VideoComposer~\cite{wang2023videocomposer} & 50.0 & 94.2 & \textbf{70.0} & 59.4 & 96.2 & 89.8 & 76.6 \\
& ControlVideo~\cite{zhao2023controlvideo} & \textbf{63.3} & 95.0 & 10.0 & 65.1 & 96.5 & 92.3 & 70.4 \\
& VACE & 56.7 & \textbf{96.1} & 60.0 & 66.4 & 98.8 & \textbf{94.1} & 78.7 \\
& Ours & 62.9 & 95.1 & 68.4 & \textbf{70.4} & \textbf{99.0} & 94.1 & \textbf{81.6} \\
\midrule
\multirow{5}{*}{Pose} 
& Text2Video-Zero~\cite{khachatryan2023text2videozero} & 57.6 & 87.7 & \textbf{100.0} & \textbf{70.7} & 79.7 & 84.8 & 80.1 \\
& ControlVideo~\cite{zhao2023controlvideo} & \textbf{65.4} & 94.6 & 25.0 & 65.3 & 97.3 & 92.8 & 73.4 \\
& Follow-Your-Pose~\cite{ma2024followyourpose} & 48.8 & 86.8 & \textbf{100.0} & 67.4 & 90.1 & 80.2 & 78.9 \\
& VACE & 60.2 & \textbf{94.9} & 75.0 & 64.7 & 98.6 & \textbf{94.8} & 81.4 \\
& Ours & 63.6 & 93.1 & 84.2 & 68.5 & \textbf{98.9} & 94.0 & \textbf{83.7} \\
\midrule
\multirow{3}{*}{Inpainting} 
& ProPainter~\cite{zhou2023propainter} & 44.7 & 95.6 & 50.0 & 61.6 & 99.0 & 93.0 & 74.0 \\
& VACE & 51.3 & \textbf{96.3} & 50.0 & 60.4 & 99.1 & 94.6 & 75.3 \\
& Ours & \textbf{59.7} & 96.3 & \textbf{55.0} & \textbf{68.8} & \textbf{99.3} & \textbf{95.4} & \textbf{79.1} \\
\midrule
\multirow{4}{*}{Outpainting} 
& Follow-Your-Canvas~\cite{chen2024followyourcanvas} & 53.3 & 96.0 & 5.0 & \textbf{69.5} & 98.1 & \textbf{95.4} & 69.5 \\
& M3DDM~\cite{fan2023m3ddm} & 53.3 & 95.9 & 30.0 & 65.1 & 99.2 & 93.7 & 72.9 \\
& VACE & \textbf{57.0} & 96.6 & 30.0 & 69.5 & 99.2 & 94.5 & 74.5 \\
& Ours & 56.1 & \textbf{96.7} & \textbf{45.0} & 68.3 & \textbf{99.4} & 95.4 & \textbf{76.8} \\
\bottomrule
\end{tabular}
\vspace{1mm}
\small

AQ: Aesthetic Quality, BC: Background Consistency, DD: Dynamic Degree, IQ: Imaging Quality, MS: Motion Smoothness, SC: Subject Consistency.
\end{table*}


Table~\ref{tab:vbench_comparison} reports VBench metrics on the VACE Benchmark. Our method achieves the highest average score on all four tasks.

\paragraph{Depth and pose.} Our method outperforms VACE by 2.9 points on depth and 2.3 on pose, while maintaining high dynamic degree (68.4 depth, 84.2 pose) and avoiding the over-constraining failure mode of methods like ControlVideo (DD of 10--25).

\paragraph{Inpainting and outpainting.} We use the same inpainting LoRA for both tasks. Our method outperforms VACE by 3.8~points on inpainting and 2.3~points on outpainting, with gains driven by substantially higher aesthetic quality (+8.4) and imaging quality (+8.4) on inpainting.

\begin{figure*}[t]
\centering
\setlength{\tabcolsep}{0.5pt}
\renewcommand{\arraystretch}{0.3}
\newcommand{\frameimg}[1]{\includegraphics[width=0.08\linewidth]{figures/images/qual_comp/#1}}
\small
\begin{tabular}{@{} *{4}{c} @{\hskip 4pt} *{4}{c} @{\hskip 4pt} *{4}{c} @{}}
\multicolumn{4}{c}{Control Input}
& \multicolumn{4}{c}{\textcolor{blue}{Ours}}
& \multicolumn{4}{c}{\textcolor{red}{VACE}} \\[2pt]
%
\frameimg{000050_control_0} & \frameimg{000050_control_1} & \frameimg{000050_control_2} & \frameimg{000050_control_3} & %\frameimg{000050_control_4} &
\frameimg{000050_ours_0} & \frameimg{000050_ours_1} & \frameimg{000050_ours_2} & \frameimg{000050_ours_3} & %\frameimg{000050_ours_4} &
\frameimg{000050_vace_0} & \frameimg{000050_vace_1} & \frameimg{000050_vace_2} & \frameimg{000050_vace_3} \\ %\frameimg{000050_vace_4} \\
%
\frameimg{000051_control_0} & \frameimg{000051_control_1} & \frameimg{000051_control_2} & \frameimg{000051_control_3} & %\frameimg{000051_control_4} &
\frameimg{000051_ours_0} & \frameimg{000051_ours_1} & \frameimg{000051_ours_2} & \frameimg{000051_ours_3} & %\frameimg{000051_ours_4} &
\frameimg{000051_vace_0} & \frameimg{000051_vace_1} & \frameimg{000051_vace_2} & \frameimg{000051_vace_3} \\ %\frameimg{000051_vace_4} \\
%
\frameimg{000052_control_0} & \frameimg{000052_control_1} & \frameimg{000052_control_2} & \frameimg{000052_control_3} & %\frameimg{000052_control_4} &
\frameimg{000052_ours_0} & \frameimg{000052_ours_1} & \frameimg{000052_ours_2} & \frameimg{000052_ours_3} & %\frameimg{000052_ours_4} &
\frameimg{000052_vace_0} & \frameimg{000052_vace_1} & \frameimg{000052_vace_2} & \frameimg{000052_vace_3} \\[4pt] %\frameimg{000052_vace_4} \\[4pt]
%
\frameimg{000076_control_0} & \frameimg{000076_control_1} & \frameimg{000076_control_2} & \frameimg{000076_control_3} & %\frameimg{000076_control_4} &
\frameimg{000076_ours_0} & \frameimg{000076_ours_1} & \frameimg{000076_ours_2} & \frameimg{000076_ours_3} & %\frameimg{000076_ours_4} &
\frameimg{000076_vace_0} & \frameimg{000076_vace_1} & \frameimg{000076_vace_2} & \frameimg{000076_vace_3} \\ %\frameimg{000076_vace_4} \\
%
\frameimg{000071_control_0} & \frameimg{000071_control_1} & \frameimg{000071_control_2} & \frameimg{000071_control_3} & %\frameimg{000071_control_4} &
\frameimg{000071_ours_0} & \frameimg{000071_ours_1} & \frameimg{000071_ours_2} & \frameimg{000071_ours_3} & %\frameimg{000071_ours_4} &
\frameimg{000071_vace_0} & \frameimg{000071_vace_1} & \frameimg{000071_vace_2} & \frameimg{000071_vace_3} \\ %\frameimg{000071_vace_4} \\
%
\frameimg{000073_control_0} & \frameimg{000073_control_1} & \frameimg{000073_control_2} & \frameimg{000073_control_3} & %\frameimg{000073_control_4} &
\frameimg{000073_ours_0} & \frameimg{000073_ours_1} & \frameimg{000073_ours_2} & \frameimg{000073_ours_3} & %\frameimg{000073_ours_4} &
\frameimg{000073_vace_0} & \frameimg{000073_vace_1} & \frameimg{000073_vace_2} & \frameimg{000073_vace_3} \\ %\frameimg{000073_vace_4} \\
\end{tabular}
\caption{Qualitative comparison on the VACE Benchmark (depth and pose). Each triplet: control input, \textcolor{blue}{ours}, \textcolor{red}{VACE}~\cite{jiang2025vace}. Our outputs show higher structural fidelity, consistent with Table~\ref{tab:vbench_comparison}.}
\label{fig:qualitative_comparison}
\end{figure*}


Figure~\ref{fig:qualitative_comparison} presents a qualitative comparison with VACE on the benchmark. Our outputs exhibit higher structural fidelity while maintaining natural motion and visual quality.

\subsection{Extended Modalities}
\label{sec:modalities}

\begin{figure*}[t]
\centering
\definecolor{oursblue}{HTML}{3B9DD9}%
\newlength{\mgfw}\setlength{\mgfw}{0.184\textwidth}%
\newcommand{\mgimg}[1]{\includegraphics[width=\mgfw]{figures/images/modality_gallery/#1}}%
\newcommand{\mgout}[1]{{\setlength{\fboxsep}{0pt}\setlength{\fboxrule}{0.5pt}%
  \fcolorbox{oursblue}{white}{\includegraphics[width=\dimexpr\mgfw-1pt\relax]{figures/images/modality_gallery/#1}}}}%
\newcommand{\modpair}[2]{%
  \begin{minipage}[c]{0.04\textwidth}%
    \centering\rotatebox[origin=c]{90}{\scriptsize\textbf{#1}}%
  \end{minipage}%
  \hfill
  \begin{minipage}[c]{0.95\textwidth}%
    \mgimg{#2_ctrl_1}\hfill\mgimg{#2_ctrl_2}\hfill\mgimg{#2_ctrl_3}\hfill\mgimg{#2_ctrl_4}\hfill\mgimg{#2_ctrl_5}\\[0.5pt]%
    \mgout{#2_out_1}\hfill\mgout{#2_out_2}\hfill\mgout{#2_out_3}\hfill\mgout{#2_out_4}\hfill\mgout{#2_out_5}%
  \end{minipage}%
}%
%
\modpair{Canny}{canny}\\[3pt]
\modpair{Pose}{pose}\\[3pt]
\modpair{Depth}{depth}\\[3pt]
\modpair{Camera (Img)}{camera_image}\\[3pt]
\modpair{Camera (Vid)}{camera_video}\\[3pt]
\modpair{Sparse Track.}{sparse_tracking}\\[3pt]
\modpair{Cut-on-Action}{cut_on_action}\\[3pt]
\modpair{Inpainting}{inpainting}%
%
\caption{Partial gallery of control modalities. Each row pair shows control input (top) and generated output (bottom, \textcolor{oursblue}{blue border}) across five sampled frames. Each modality is an independent LoRA trained in 200--15{,}000 steps. We provide additional video examples in the supplementary.}
\label{fig:modality_gallery}
\end{figure*}


Beyond the benchmark controls, we demonstrate the breadth of modalities the framework supports (Figure~\ref{fig:modality_gallery}). Each modality uses its own LoRA, and adding a new one requires no retraining of existing ones. We support ControlNet-style modalities (depth, pose, canny), video editing (inpainting/outpainting, detailing), and composited controls (e.g., masked depth with pose for Blender rendering). We also train a \emph{sparse tracks} LoRA for point-trajectory-based motion control, similar to ATI~\cite{wang2025ati}, with tracks extracted via AllTracker~\cite{harley2025alltracker}. More results are demonstrated in the supplementary video.

\paragraph{Camera trajectory control.}
\label{sec:camera}
We support two modes of camera control: (1)~generating diverse camera motions from a single input image, and (2)~re-rendering an existing video at a new camera trajectory while preserving the original scene motion. For the latter, we estimate full camera parameters (extrinsics and intrinsics, including per-frame FOV) from the source video using SpatialTrackerV2~\cite{xiao2025spatialtracker}, then re-render each frame at the desired camera configuration. Optionally, rendering from a different timestamp retimes the output. Training uses a synthetic dataset of synchronized moving cameras, similar to ReCamMaster~\cite{bai2025recammaster}.
Unlike ReCamMaster~\cite{bai2025recammaster}, which controls only camera extrinsics, our camera LoRAs also control intrinsics, specifically field of view (FOV). This enables simulating focal length changes and effects such as the dolly zoom (``vertigo effect''), which is impossible with extrinsics-only methods.

\begin{table}[t]
\centering
\caption{Camera control on the ReCamMaster Benchmark (200 videos). Baselines from~\protect\cite{bai2025recammaster}. $^\dagger$COLMAP Structure-from-Motion (SfM) fails on 27\% of our videos; SpatialTrackerV2~\protect\cite{xiao2025spatialtracker} (all videos) yields 3.55$^\circ$.}
\label{tab:camera_comparison}
\setlength{\tabcolsep}{6pt}
\begin{tabular}{lcc}
\toprule
Method & CLIP-F (\%) $\uparrow$ & RotErr ($^\circ$)  \\
\midrule
GCD~\cite{van2024gcd} & 95.66 & 2.27 \\
Traj-Attn~\cite{xiao2025trajattn} & 96.52 & 2.18 \\
DaS~\cite{cheng2025das} & 98.32 & 1.45 \\
ReCamMaster~\cite{bai2025recammaster} & 98.74 & \textbf{1.22} \\
\midrule
\textbf{Ours} & \textbf{99.13} & 6.00$^\dagger$ \\
\bottomrule
\end{tabular}
\end{table}


Table~\ref{tab:camera_comparison} compares our camera control against dedicated methods on the ReCamMaster Benchmark~\cite{bai2025recammaster}. We evaluate on 200~randomly sampled videos across 10~trajectory types. Our method achieves the highest CLIP-F score (99.13\%), surpassing ReCamMaster (98.74\%). Our COLMAP-based RotErr is 6.00$^\circ$, though SfM fails on 27\% of our videos, likely underestimating the true average error. SpatialTrackerV2 tracking on all 200~videos yields 3.55$^\circ$ (not directly comparable to baselines). While dedicated camera control architectures achieve lower rotation error, our camera LoRA is a lightweight adapter on a general-purpose audio-visual model and additionally controls camera intrinsics (FOV), which extrinsics-only methods cannot.

\paragraph{Audio-visual applications.}
\label{sec:av}
Audio-visual modalities are qualitatively different from the video-only controls above: training uses audio pairs, and the generated output spans both modalities (see supplementary video for qualitative results). We demonstrate two audio modalities: \emph{audio intensity control} and \emph{speech-to-ambient}, plus the cross-modal \emph{who-is-talking} control. Each audio LoRA is trained on audio-only pairs, yet at inference time generates both audio and video in a single joint pass: training on a single modality, deploying on both.

\emph{Audio intensity control} generates audio whose temporal dynamics follow visual content, addressing the same task as ReWaS~\cite{jeong2024rewas} and CAFA~\cite{benita2025cafa}. We train a single LoRA on the joint model rather than a dedicated video encoder and adapter on a unimodal backbone. The reference canvas carries the original video with its audio replaced by an energy-envelope sonification.

\begin{table}[t]
\centering
\caption{Audio intensity evaluation on VGGSound. Baselines are dedicated V2A methods; ours generates both audio and video jointly via a single LoRA. Baseline numbers from~\protect\cite{benita2025cafa}. \textbf{Bold}: best per metric.}
\label{tab:audio_intensity}
\setlength{\tabcolsep}{3.5pt}
\begin{tabular}{lccccc}
\toprule
Method & FAD & KL & IS$\uparrow$ & IB$\uparrow$ & Train Data \\
\midrule
ReWaS~\cite{jeong2024rewas}      & 14.71 & 2.69 & 8.45  & 0.15 & 160K \\
CAFA~\cite{benita2025cafa}       & 12.60 & 2.02 & 13.45 & 0.21 & 200K \\
MMAudio~\cite{cheng2025mmaudio}  & \textbf{5.32}  & \textbf{1.64} & 17.18 & \textbf{0.33} & 180K \\
\midrule
Ours & 57.25 & 7.74 & \textbf{34.51} & 0.13 & 7.8K \\
\bottomrule
\end{tabular}
\vspace{1mm}

\small
FAD: Fr\'{e}chet Audio Distance (PANNS), KL: KL divergence (PANNS), IS: Inception Score (PANNS), IB: ImageBind audio-visual similarity. ReWaS and CAFA numbers use the TV2A configuration from~\cite{benita2025cafa} Table~2; MMAudio uses the V2A configuration (video input only, no text).
\end{table}


Table~\ref{tab:audio_intensity} compares our intensity LoRA against dedicated video-to-audio methods on 254 VGGSound~\cite{chen2020vggsound} test samples. We report FAD, KL, IS (via PANNS~\cite{kong2020panns}) and audio-visual alignment (IB via ImageBind~\cite{girdhar2023imagebind}), following~\cite{benita2025cafa}. Our method achieves the highest IS (34.51), indicating diverse and class-distinctive audio generation. As expected, FAD and KL are higher because our model generates both audio and video jointly in a single pass rather than specializing in audio-only synthesis, and trains on 20--25$\times$ less data (7.8K vs.\ 160--200K samples). Notably, our IB score (0.13) is within 0.02 of ReWaS (0.15), a dedicated V2A model, despite our fundamentally different architecture: a single lightweight LoRA on a joint audio-visual backbone with no dedicated video encoder. Additionally, our audio quality (FAD) is capped by the quality of the LTX-2 vocoder.

\emph{Speech-to-ambient} embeds clean speech within ambient sounds matching a text-described scene. It operates in two modes: \emph{audio replacement} (AV$\to$A), where the video's audio is replaced while preserving the video, and \emph{joint generation} (A$\to$AV), where both audio and video are generated from the speech signal and a text prompt. The LoRA is trained on 2{,}600 paired samples in 5{,}000 steps.

\emph{Who-is-talking control} demonstrates an inherently cross-modal control: given a reference encoding spatial layout (bounding boxes) and temporal activity (which speaker is active when), the model generates multi-person talking video with synchronized lip motion and audio. The reference is abstract (colored rectangles on a black background, colored when speaking and gray when silent), yet the parallel canvas maps this to realistic video with joint audio.

To compare with MultiTalk~\cite{kong2025multitalk}, we start from the \emph{same raw inputs}: per-speaker identity references and separated audio tracks. We generate a first frame conditioned on the reference images, derive a bounding-box activity map via voice activity detection, and mix all tracks into a single audio signal. Our method supports an arbitrary number of speakers, unlike MultiTalk which is limited to two, and can generate audio jointly when none is provided.

\begin{table}[t]
\centering
\caption{Talking-head comparison on HDTF test set. Sync-C (confidence, $\uparrow$) and Sync-D (distance, ) measure lip-audio synchronization via SyncNet. E-FID measures expression quality. FID measures visual quality.
Baseline numbers are from~\cite{kong2025multitalk}.}
\label{tab:whoistalking}
\setlength{\tabcolsep}{5pt}
\begin{tabular}{lcccc}
\toprule
Method & Sync-C $\uparrow$ & Sync-D  & E-FID  & FID  \\
\midrule
AniPortrait~\cite{wei2024aniportrait} & 3.09 & 10.94 & 1.32 & 32.83 \\
Hallo3~\cite{cui2025hallo3} & 6.55 & 8.49 & 1.12 & 33.98 \\
Sonic~\cite{ji2025sonic} & 8.35 & 6.43 & 1.22 & 29.53 \\
MultiTalk~\cite{kong2025multitalk} & \textbf{8.54} & \textbf{6.69} & 1.00 & 24.01 \\
\midrule
\textbf{Ours} & 4.50 & 10.31 & \textbf{0.18} & \textbf{12.31} \\
\bottomrule
\end{tabular}
\end{table}


Table~\ref{tab:whoistalking} compares on the HDTF benchmark. Our method achieves strong expression quality (E-FID~0.18) and visual fidelity (FID~12.31), outperforming all baselines on both. Lip-sync scores lag behind dedicated methods, as expected for a single general-purpose LoRA without purpose-built audio cross-attention.

\subsection{Ablations and Analysis}
\label{sec:efficiency}

Individual LoRAs range from 200~steps to 15{,}000~steps, with an aggregate budget of ${\sim}$55{,}000 steps across all trained modalities, less than one third of VACE's~\cite{jiang2025vace} 200{,}000 and comparable to a single specialized method like BulletTime~\cite{wang2025bullettime} (40K iterations). See Figure~\ref{fig:training_efficiency} in the supplementary for a visual comparison.

\paragraph{Training convergence.}
Training loss plateaus early for spatially-aligned controls. For depth, VBench scores reach 81.1 at 1{,}000 steps and 81.6 at 3{,}000 steps. The rapid convergence is consistent with LoRA fine-tuning existing weights rather than training new projections from scratch.

\paragraph{LoRA rank.}
LoRA ranks 32, 64, and 128 for depth yield VBench averages of 80.9, 81.3, and 81.6, a spread of less than 1~point. We use rank~128 as the default for a small quality margin.

\paragraph{Parallel canvas vs.\ spatial concatenation.}
A concatenation-based baseline (analogous to IC-LoRA~\cite{huang2024iclora}) trained with the same data, schedule, and rank fails to follow the depth signal (Figure~\ref{fig:canvas_vs_concat}), confirming the conditioning mechanism is the critical factor.

\paragraph{Generalization from synthetic data.}
Several LoRAs (camera-from-video, cut-on-action, local edit) are trained exclusively on synthetic or generated data yet generalize to real-world videos without fine-tuning. See the supplementary for details.
